\documentclass[11pt]{article}

\usepackage[margin=2.5cm]{geometry}
\usepackage{amsmath,amssymb,amsthm,mathtools}
\usepackage{enumitem}
\usepackage{xcolor}
\usepackage{hyperref}

\newtheorem{theorem}{Theorem}[section]
\newtheorem{proposition}[theorem]{Proposition}
\newtheorem{lemma}[theorem]{Lemma}
\newtheorem{corollary}[theorem]{Corollary}
\theoremstyle{definition}
\newtheorem{definition}[theorem]{Definition}
\newtheorem{remark}[theorem]{Remark}

\newcommand{\One}{\textsc{One}}
\newcommand{\All}{\textsc{All}}
\newcommand{\N}{\mathbb{N}}
\newcommand{\R}{\mathbb{R}}
\newcommand{\Z}{\mathbb{Z}}
\newcommand{\E}{\mathbb{E}}

%% ---------------------------------------------------------------------
%%  DRAFT MARKER.  \gap{label}{text} flags a step that is NOT yet
%%  complete.  It renders as a conspicuous red box in the PDF.  Every
%%  such box marks a place where work must continue; the manuscript is
%%  written in finished form around them on purpose, so that the
%%  narrative and the precise proof obligations are visible.
%% ---------------------------------------------------------------------
\newcommand{\gap}[2]{%
  \par\smallskip\noindent
  \fcolorbox{red!70!black}{red!7}{%
    \parbox{\dimexpr\linewidth-2\fboxsep-2\fboxrule}{%
      {\bfseries\color{red!70!black}$\blacktriangleright$~OPEN GAP --- #1.
       Work continues here.}\\[2pt]#2}}%
  \par\smallskip}

\title{Log-periodic oscillation in the all-heads coin game}
\author{Peter Pfaffelhuber}
\date{\today}

\begin{document}
\maketitle

\begin{abstract}
  We study the large-$n$ behaviour of the all-heads coin game below
  the fair coin. Two phenomena are established. First, for every
  $p\in(0,\tfrac12)$ the value $b_{n,p}$ of the greedy strategy
  \All{} is asymptotically \emph{log-periodic}: it is asymptotic to a
  continuous, non-constant function of $\log n$ with period
  $\log\bigl(1/(1-p)\bigr)$, and in particular does not converge.
  Second --- and this is the main result --- the \emph{optimal}
  winning probability $w_{n,p}$ inherits this oscillation: for every
  $p\in(0,\tfrac12)$ the sequence $n\mapsto w_{n,p}$ is not eventually
  monotone and has infinitely many strict local maxima and minima.
  The proof reduces the statement to two one-sided assertions, treats
  the increasing side by an elementary Bellman estimate, and the
  decreasing side by an exact transfer to strategy~\All{} followed by
  a renormalisation argument that propagates the oscillation across
  geometric scales. This paper is the asymptotic companion to the
  exact and perturbative analysis of~\cite{Paper1}.
\end{abstract}

\noindent\emph{\small Working draft. The one step still to be
completed is flagged in the text by a boxed}
{\small\bfseries\color{red!70!black}OPEN GAP} \emph{\small marker
(Gap~D); Section~\ref{sec:status} consolidates its status.}

%% =====================================================================
\section{Introduction}\label{sec:intro}
%% =====================================================================

\subsection{The game}

Fix $p\in(0,1)$ and write $q:=1-p$. A player holds $n$ coins, each of
which lands heads independently with probability~$p$. In each round
all remaining coins are tossed; the player must set aside at least one
coin showing heads, and loses if no coin shows heads. The game is won
once every coin has been set aside. Writing $w_{n,p}$ for the optimal
winning probability, the dynamic-programming principle gives the
Bellman equation
\begin{equation}\label{eq:bellman}
  w_{n,p}=p^{n}+\sum_{j=1}^{n-1}\binom{n}{j}p^{\,n-j}q^{\,j}
          \max_{j\le m\le n-1}w_{m,p},\qquad w_{0,p}:=1 .
\end{equation}
Two strategies are natural: \One{} sets aside a single head each round
(value $a_{n,p}$), and \All{} sets aside \emph{every} head. The
winning probability $b_{n,p}$ of strategy~\All{} is defined by the
linear recursion
\begin{equation}\label{eq:brec}
  b_{n,p}=\sum_{j=0}^{n-1}\binom{n}{j}p^{\,n-j}q^{\,j}\,b_{j,p},
  \qquad b_{0,p}:=1,
\end{equation}
obtained by conditioning on the number $j$ of coins showing tails in
the first round: those $j$ coins are re-tossed, while the $n-j\ge1$
heads are set aside for good (the term $j=0$ has weight $p^{n}$ and
wins immediately). The companion paper~\cite{Paper1} determines
$w_{n,p}$
\emph{exactly} at $p=\tfrac12$, identifies the limit
$W(p)=\lim_n w_{n,p}$ for $p>\tfrac12$, and gives a first-order
perturbation expansion for $p<\tfrac12$ near $\tfrac12$. All of those
results are machine-checked in Lean~4.

The present paper concerns the regime left open by a perturbative
analysis: the genuine large-$n$ behaviour for $p<\tfrac12$, which
turns out to be \emph{oscillatory} rather than convergent.

\subsection{Main results}\label{sec:main}

We state the two theorems at once; everything else in the paper is in
their service.

\begin{theorem}[oscillation of the optimal value]\label{thm:main}
  For every $p\in(0,\tfrac12)$ the sequence $n\mapsto w_{n,p}$ is not
  eventually monotone. It has infinitely many strict local maxima and
  infinitely many strict local minima.
\end{theorem}

\begin{theorem}[log-periodicity of strategy \All]\label{thm:ALL}
  For every $p\in(0,\tfrac12)$ there is a continuous, strictly
  positive, non-constant function $\Phi_p:\mathbb R\to\mathbb R$,
  periodic with period $L:=\log(1/q)$, such that
  \[
    b_{n,p}=\Phi_p(\log n)\,\bigl(1+o(1)\bigr)\qquad(n\to\infty).
  \]
  In particular $b_{n,p}$ does not converge:
  $\liminf_n b_{n,p}<\limsup_n b_{n,p}$.
\end{theorem}

Theorem~\ref{thm:main} sharpens, and corrects the natural first
guess about, the first-order picture of~\cite{Paper1}: that analysis
shows $w_{n,p}$ to be eventually increasing \emph{at first order in}
$\tfrac12-p$, but the exact sequence is not eventually monotone at
all. The oscillation is non-perturbative --- its amplitude is
exponentially small in $1/(\tfrac12-p)$ --- which is why it is
invisible to any finite-order expansion.

\subsection{Strategy of proof}\label{sec:strategy}

The proof of Theorem~\ref{thm:main} has a rigid skeleton, which we
record now so that the role of each later section is clear. For a
real sequence $(x_n)$ put $d_n:=x_{n+1}-x_n$.

\begin{quote}
\textbf{Step 0 (reduction).} A bounded sequence has infinitely many
strict local maxima and minima as soon as it is \emph{not eventually
monotone}, i.e.\ as soon as $d_n>0$ infinitely often \emph{and}
$d_n<0$ infinitely often (Lemma~\ref{lem:reduction}). Thus
Theorem~\ref{thm:main} splits into
\[
  \text{(A) }(w_{n,p})\text{ is not eventually non-decreasing,}
  \qquad
  \text{(B) }(w_{n,p})\text{ is not eventually non-increasing.}
\]

\textbf{Step A (the increasing side).} On a stretch where $w$ is
non-decreasing, the suffix-maximum in~\eqref{eq:bellman} collapses to
its right endpoint and the Bellman equation becomes the linear
strategy-\One{} recursion, whose one-step increment has the
sign-definite dominant term $-q^{n}w_{n-1}$. Since $q>p$, this forces
an eventual decrease. Statement~(A) follows
(Proposition~\ref{prop:A}); the argument is elementary and complete.

\textbf{Step B (the decreasing side).} On a stretch where $w$ is
non-increasing, the suffix-maximum collapses to its \emph{left}
endpoint, and $w$ becomes --- exactly --- a strategy-\All{} sequence
(Lemma~\ref{lem:transfer}). Strategy~\All{} oscillates
(Theorem~\ref{thm:ALL}); hence $w$ cannot have a non-increasing tail.
The transfer of the oscillation from \All{} to $w$ is carried out by
a renormalisation argument (Section~\ref{sec:renorm}): a windowed
fundamental coefficient $A(t)$ is shown to satisfy a propagation
identity across the geometric scale, from which non-convergence ---
hence non-monotonicity --- follows.
\end{quote}

The asymmetry between Steps~A and~B is structural: the increasing side
reduces to a contractive linear recursion with a sign-definite drift,
the decreasing side to a scale-invariant recursion whose solutions
oscillate. Step~A is elementary; Step~B carries the analytic weight of
the paper.

%% =====================================================================
\section{Setup and notation}\label{sec:setup}
%% =====================================================================

Throughout, $p\in(0,\tfrac12)$ is fixed, $q=1-p\in(\tfrac12,1)$, and
we use freely that $q>p$, $p/q<1$ and $2p<1$. We write
$L:=\log(1/q)>0$ and $\omega:=2\pi/L$.

The value $w_{n,p}$ satisfies~\eqref{eq:bellman}; it is a probability,
so $0<w_{n,p}\le1$, and in fact $w_{n,p}<1$ for $n\ge1$. The two
reference strategies are dominated by the optimum. Strategy~\All{} has
winning probability $b_{n,p}$, defined by the
recursion~\eqref{eq:brec}; strategy~\One{} satisfies the linear
recursion
\begin{equation}\label{eq:arec}
  a_{n,p}=p^{n}+(1-p^{n}-q^{n})\,a_{n-1,p},\qquad a_{0,p}:=1 .
\end{equation}
In both cases $w_{n,p}\ge\max(a_{n,p},b_{n,p})$, since
the optimum dominates any fixed strategy. We use the binomial
identities $\sum_{j=1}^{n-1}\binom nj p^{n-j}q^j=1-p^n-q^n$ and
$\sum_{j=0}^{n-1}\binom nj p^{n-j}q^j=1-q^n$. Where the value of $p$
is fixed we abbreviate $w_n:=w_{n,p}$, $b_n:=b_{n,p}$.

%% =====================================================================
\section{Strategy \All{}: log-periodic asymptotics}\label{sec:ALL}
%% =====================================================================

This section establishes Theorem~\ref{thm:ALL} --- the existence of
the log-periodic profile $\Phi_p$ --- save for one point, the
non-constancy of $\Phi_p$, which is Gap~D (Section~\ref{sec:base}).
The route is elementary: an exact product identity for the generating
function, followed by analytic de-Poissonisation. Neither Hayman
admissibility theory nor a Mellin contour shift is needed.

\subsection{The functional equation}

Let $B(z):=\sum_{n\ge0}b_n z^n/n!$ be the exponential generating
function of strategy~\All. The convolution structure
of~\eqref{eq:brec} yields the first-order $q$-difference equation
\begin{equation}\label{eq:Bfunc}
  B(z)=1+(e^{pz}-1)\,B(qz),
\end{equation}
and iterating it (the remainder tends to $0$ since each factor
$e^{pq^kz}-1\to0$),
\begin{equation}\label{eq:Bseries}
  B(z)=\sum_{M\ge0}\phi_M(z),\qquad
  \phi_M(z):=\prod_{k=0}^{M-1}\bigl(e^{pq^kz}-1\bigr).
\end{equation}
Probabilistically, with $X_1,\dots,X_n$ i.i.d.\ geometric,
$b_{n,p}=\Pr(\{X_1,\dots,X_n\}=\{1,\dots,\max_iX_i\})$.

\subsection{An exact product identity}\label{sec:identity}

Write $h(z):=\log(1-e^{-pz})$. From $e^{pq^kz}-1=e^{pq^kz}
(1-e^{-pq^kz})$ and $p\sum_{k<M}q^k=1-q^M$,
\begin{equation}\label{eq:phiM}
  \phi_M(z)=e^{z(1-q^M)}\,\exp S_M(z),\qquad
  S_M(z):=\sum_{k=0}^{M-1}h(q^kz).
\end{equation}
The associated \emph{convergent} harmonic sum is
\begin{equation}\label{eq:Htilde}
  \tilde H(x):=\sum_{j\ge0}h(q^{-j}x)
\end{equation}
--- here the arguments $q^{-j}x\to\infty$, so $h(q^{-j}x)\to0$
exponentially. Reindexing, $\tilde H(q^{M-1}z)=\sum_{k\le M-1}
h(q^kz)$, whence $S_M(z)=\tilde H(q^{M-1}z)-\tilde H(z/q)$.
Substituting this into~\eqref{eq:phiM}--\eqref{eq:Bseries} and setting
\begin{equation}\label{eq:Xi}
  \Xi(w):=e^{-w}\,e^{\tilde H(w/q)}
\end{equation}
yields the exact identity
\begin{equation}\label{eq:Bidentity}
  B(z)=e^{z}\,e^{-\tilde H(z/q)}\sum_{M\ge0}\Xi(q^Mz).
\end{equation}

\subsection{The profile}\label{sec:profile}

The function $\Xi$ is positive, smooth, and decays at \emph{both}
ends: $\Xi(w)\le e^{-w}$ as $w\to\infty$ (since $\tilde H\le0$), and
$\Xi(w)\to0$ faster than any power as $w\to0^+$, because
$\tilde H(x)\to-\infty$ (a sum of $\sim|\log x|/L$ terms each
$\approx\log(px)$; in fact $\tilde H(x)\sim-(\log x)^2/(2L)$). Hence
$\sum_{M\ge0}\Xi(q^Mz)$ is a harmonic sum of a bump. Its bi-infinite
completion
\begin{equation}\label{eq:Phidef}
  \Phi(t):=\sum_{M\in\mathbb Z}\Xi(q^Me^{t})
\end{equation}
converges locally uniformly; $\Phi$ is continuous, strictly positive,
bounded, and \emph{exactly} $L$-periodic. The omitted terms $M\le-1$
are $O(e^{-cz})$ and $\tilde H(z/q)=O(e^{-cz})$,
so~\eqref{eq:Bidentity} gives, on the positive axis,
\begin{equation}\label{eq:Basymp}
  B(z)=e^{z}\,\Phi(\log z)\,\bigl(1+O(e^{-cz})\bigr),
  \qquad z\to+\infty .
\end{equation}
The Fourier coefficients of $\Phi$ are $\widehat\Phi(m)=
L^{-1}\,\mathcal M[\Xi](2\pi i m/L)$, with $\mathcal M[\Xi]$ entire
($\Xi$ decays faster than any power at both ends).

\subsection{Transfer to coefficients}\label{sec:transfer-coeff}

The passage from $B(z)$ to $b_n=n!\,[z^n]B(z)$ is analytic
de-Poissonisation~\cite{JS}. Write $\Psi(z):=B(z)e^{-z}$; the required
sector control is read off the product~\eqref{eq:Bseries}.

\begin{lemma}[sector control]\label{lem:sector}
  There are constants $C,\gamma$, depending only on $p$, such that for
  $z=re^{i\theta}$ with $r$ large,
  \[
    |B(re^{i\theta})|\le
      \begin{cases}
        C\,r^{\gamma}\,e^{r\cos\theta}, & |\theta|\le\pi/2,\\[2pt]
        C\,r^{\gamma}, & \pi/2\le|\theta|\le\pi .
      \end{cases}
  \]
\end{lemma}

\begin{proof}
  Let $k^\ast$ be least with $pq^{k^\ast}r\le1$; then $k^\ast\le
  1+\log(pr)/L$. For $k<k^\ast$ use $|e^{pq^kz}-1|\le1+e^{pq^kr
  \cos\theta}$; for $k\ge k^\ast$ use $|e^{w}-1|\le|w|\,e^{(\Re w)^+}$,
  i.e.\ $|e^{pq^kz}-1|\le pq^kr\,e^{(pq^kr\cos\theta)^+}$. If
  $\cos\theta\ge0$, the early factors give $\prod_{k<k^\ast}2\,
  e^{pq^kr\cos\theta}\le2^{k^\ast}e^{r\cos\theta}$ (as
  $\sum_kpq^k=1$), and the late factors $\prod_{k\ge k^\ast}e\,pq^kr$
  decay super-exponentially in $M$; summing over $M$ a geometric-type
  series gives the first bound, with $\gamma=\log2/L$. If
  $\cos\theta\le0$, the early factors are each $\le2$ and the late
  ones $\le pq^kr$, giving the second bound.
\end{proof}

Lemma~\ref{lem:sector} furnishes both de-Poissonisation hypotheses.
Inside the cone $|\arg z|\le\pi/2$, $|\Psi(z)|=|B(z)|\,
e^{-r\cos\theta}\le Cr^{\gamma}$ is polynomially bounded; outside any
cone $|\arg z|\ge\delta$, $|B(z)|\le Cr^{\gamma}e^{r\cos\delta}\le
e^{\alpha r}$ with $\alpha=\tfrac12(1+\cos\delta)<1$.
De-Poissonisation then gives $b_n=\Psi(n)\,(1+O(1/n))$,
and~\eqref{eq:Basymp} yields $\Psi(n)=\Phi(\log n)(1+O(e^{-cn}))$.
Hence
\begin{equation}\label{eq:bnasymp}
  b_{n,p}=\Phi_p(\log n)\,\bigl(1+O(1/n)\bigr),
\end{equation}
with $\Phi_p:=\Phi$ continuous, strictly positive and $L$-periodic.
This is Theorem~\ref{thm:ALL} apart from the non-constancy of
$\Phi_p$.

\subsection{Non-constancy}\label{sec:nonconst}

By~\eqref{eq:bnasymp}, $b_{n,p}$ converges if and only if $\Phi_p$ is
constant, if and only if $\mathcal M[\Xi](2\pi im/L)=0$ for every
$m\ne0$. At $p=\tfrac12$ one has $b_n\equiv\tfrac12$ and
$\Phi_p\equiv\tfrac12$. For $p<\tfrac12$ non-constancy is expected but
unproved; it is implied by Gap~D (Section~\ref{sec:base}) and is the
arithmetic heart of the paper.

\begin{remark}
  The harmonic sum $\tilde H$ \emph{does} genuinely oscillate: its
  Mellin transform $-\Gamma(s)p^{-s}\zeta(s+1)/(1-q^{s})$
  (fundamental strip $\Re s>0$) has simple poles at $s_m=2\pi im/L$
  with residues $-L^{-1}\Gamma(s_m)p^{-s_m}\zeta(s_m+1)$, non-zero by
  the non-vanishing of $\zeta$ on $\Re s=1$. This is why $\Phi_p$
  \emph{should} be non-constant; but $\Phi_p$'s Fourier coefficients
  $L^{-1}\mathcal M[\Xi](s_m)$ depend on $\Xi=e^{-w}e^{\tilde H(w/q)}$
  non-linearly, and the non-vanishing does not transfer formally.
\end{remark}

%% =====================================================================
\section{Step A: \texorpdfstring{$w_{n,p}$}{w(n,p)} is not eventually
  increasing}\label{sec:A}
%% =====================================================================

\begin{proposition}\label{prop:A}
  For $p\in(0,\tfrac12)$ there is no $M$ with $w_n\ge w_{n-1}$ for all
  $n>M$. That is, $w_n<w_{n-1}$ for infinitely many~$n$.
\end{proposition}

\begin{proof}
  Suppose $w_n\ge w_{n-1}$ for all $n>M$, i.e.\ $w$ is non-decreasing
  on $\{M,M+1,\dots\}$. Fix $n>M+1$ and split the Bellman
  sum~\eqref{eq:bellman} at $j=M$, writing $S_{j,n}:=\max_{j\le m\le
  n-1}w_m$. For $j\ge M$ the index window lies in the
  non-decreasing range, so $S_{j,n}=w_{n-1}$; for every $j$ one has
  $w_{n-1}\le S_{j,n}\le1$. With $U_n:=\sum_{j=1}^{M-1}\binom nj
  p^{n-j}q^j$ and $T_n:=\sum_{j=1}^{M-1}\binom nj p^{n-j}q^j S_{j,n}$,
  \[
    w_n=p^n+T_n+(1-p^n-q^n-U_n)\,w_{n-1},
  \]
  hence
  \[
    w_n-w_{n-1}=p^n(1-w_{n-1})-q^n w_{n-1}+(T_n-U_n w_{n-1}).
  \]
  Since $w_{n-1}\le S_{j,n}\le1$ one has $0\le T_n-U_nw_{n-1}\le U_n$,
  and $U_n\le(M-1)\,n^{M-1}p^{\,n-M+1}=C_1\,n^{M-1}p^n$ with
  $C_1=(M-1)p^{1-M}$. With $c:=w_M>0$ and $w_{n-1}\ge c$,
  \[
    w_n-w_{n-1}\le\bigl(1+C_1n^{M-1}\bigr)p^n-c\,q^n .
  \]
  As $p/q<1$, the ratio $(1+C_1n^{M-1})(p/q)^n\to0$, so the
  right-hand side is negative for all large~$n$ --- contradicting
  $w_n\ge w_{n-1}$.
\end{proof}

Proposition~\ref{prop:A} is statement~(A) of Section~\ref{sec:strategy}
and is fully rigorous.

%% =====================================================================
\section{Step B, part 1: reduction and the transfer
  lemma}\label{sec:transfer}
%% =====================================================================

\begin{lemma}[reduction]\label{lem:reduction}
  A bounded real sequence $(x_n)$ that is neither eventually
  non-decreasing nor eventually non-increasing has infinitely many
  strict local maxima and infinitely many strict local minima,
  provided $x_n\ne x_{n+1}$ for all large~$n$.
\end{lemma}

\begin{proof}
  Not eventually non-decreasing means $d_n:=x_{n+1}-x_n<0$ infinitely
  often; not eventually non-increasing means $d_n>0$ infinitely
  often. Between a sign change of $d$ from $-$ to $+$ lies a local
  minimum, between $+$ and $-$ a local maximum; infinitely many sign
  changes yield infinitely many of each, strict under $x_n\ne x_{n+1}$.
\end{proof}

\begin{lemma}[transfer to strategy \All]\label{lem:transfer}
  Suppose $w_n\ge w_{n+1}$ for all $n\ge M$. Define $x_0:=1$,
  $x_j:=\max_{j\le m\le M}w_m$ for $1\le j\le M-1$, and $x_j:=w_j$ for
  $j\ge M$. Then $x_n=w_n$ for $n\ge M$, and $(x_n)$ satisfies the
  exact strategy-\All{} recursion
  \[
    x_n=\sum_{j=0}^{n-1}\binom nj p^{\,n-j}q^{\,j}\,x_j
    \qquad\text{for every }n>M .
  \]
\end{lemma}

\begin{proof}
  Fix $n>M$. In the Bellman equation~\eqref{eq:bellman} the
  suffix-maximum $S_{j,n}=\max_{j\le m\le n-1}w_m$ collapses to the
  \emph{left} endpoint on the non-increasing range: for $M\le j\le
  n-1$, $S_{j,n}=w_j=x_j$; for $1\le j\le M-1$,
  $S_{j,n}=\max(\max_{j\le m\le M-1}w_m,\,w_M)=\max_{j\le m\le
  M}w_m=x_j$, independent of $n$. Writing the $p^n$ term as
  $\binom n0p^nq^0x_0$ gives the displayed recursion.
\end{proof}

Lemma~\ref{lem:transfer} is the crux of Step~B: it says that a
non-increasing tail of $w$ is, exactly, a strategy-\All{} sequence
with finitely many initial values modified. The remaining task is to
show such a sequence cannot exist, by transferring the oscillation of
Theorem~\ref{thm:ALL}.

%% =====================================================================
\section{Step B, part 2: the renormalisation argument}\label{sec:renorm}
%% =====================================================================

Throughout this section $(x_n)$ is a bounded solution of the
strategy-\All{} recursion for $n>M$ --- either $b$ itself, or a
transferred sequence of Lemma~\ref{lem:transfer}.

\subsection{The recursion as a smoothing in log-scale}

With $J_n\sim\mathrm{Bin}(n,q)$ the recursion reads
$x_n=E[x_{J_n}]/(1+q^n)$. Read on the log-scale $t:=\log n$, and using
$\log(qn)=\log n-L$, one application of the recursion shifts the scale
down by exactly one period~$L$ and smooths by a kernel of width
$\sigma(t)\sim\sqrt{p/q}\;e^{-t/2}\to0$.

\subsection{The local fundamental coefficient}

Fix a window $\chi\in C_c^\infty(\R)$ with $\int\chi=1$ and
$\hat\chi(2\pi)=0$, and set
\begin{equation}\label{eq:Adef}
  A(t):=\frac1L\sum_{n>M}x_n\,e^{-i\omega\log n}\,
        \chi\!\Bigl(\tfrac{\log n-t}{L}\Bigr)\,\frac1n .
\end{equation}
This is a windowed Fourier sum at frequency $\omega$ on the
multiplicative grid, the weight $1/n$ being the log-scale element
$d(\log n)$. The condition $\hat\chi(2\pi)=0$ makes $A$ annihilate the
constant (mean-level) component, so $A$ extracts the fundamental
oscillation. If $x_n$ converges then $A(t)\to0$.

\subsection{The propagation estimate}

\begin{proposition}[propagation]\label{prop:prop}
  There are $\kappa(t)\in\mathbb C$ and a remainder $r(t)$ with
  \[
    A(t+L)=\kappa(t)\,A(t)+r(t),\qquad
    |\kappa(t)-1|=O(e^{-t}),\qquad |r(t)|=O(e^{-t}).
  \]
\end{proposition}

The factor $\kappa$ is fully under control.

\begin{lemma}[the factor $\kappa$]\label{lem:kappa}
  For $J_n\sim\mathrm{Bin}(n,q)$,
  $\kappa(n)=E[(J_n/(qn))^{i\omega}\mathbf 1_{J_n\ge1}]
  =1-\tfrac{p}{2qn}(\omega^2+i\omega)+O(n^{-3/2})$. In particular
  $|\kappa-1|=O(1/n)=O(e^{-t})$.
\end{lemma}

\begin{proof}[Proof (sketch)]
  Put $V_n=J_n/(qn)-1$; the binomial central moments give
  $E[V_n^2]=p/(qn)$ and $E[|V_n|^3]=O(n^{-3/2})$. Expanding
  $(1+V_n)^{i\omega}$ to second order on $\{|V_n|\le\tfrac12\}$ and
  bounding the tail by Hoeffding yields the claim.
\end{proof}

For the remainder we pass to the \emph{transposed} kernel. This makes
the negative-binomial structure of the recursion explicit and, as it
turns out, removes the need for any limit theorem. Write
$a_n(t):=n^{-1}e^{-i\omega\log n}\chi\bigl(\tfrac{\log n-t}L\bigr)$, so
that $A(t)=\frac1L\sum_{n>M}x_n\,a_n(t)$.

\begin{lemma}[transposed kernel and the negative binomial]%
\label{lem:transposed}
  Substituting the recursion
  $x_n=(1+q^n)^{-1}\sum_{j=0}^{n}P(\mathrm{Bin}(n,q)=j)\,x_j$
  into $A(t+L)$ and exchanging the (finite) order of summation,
  \[
    A(t+L)=\frac1L\sum_{j}x_j\,W(j,t),\qquad
    W(j,t)=\sum_{n\ge j}\frac{P(\mathrm{Bin}(n,q)=j)}{1+q^n}\,
           a_n(t+L).
  \]
  The transposed weight has \emph{exact} total mass
  \[
    \sum_{n\ge j}P(\mathrm{Bin}(n,q)=j)
      =q^{j}\sum_{k\ge0}\binom{j+k}{j}p^{k}
      =q^{j}(1-p)^{-(j+1)}=\frac1q ,
  \]
  and $q\,P(\mathrm{Bin}(n,q)=j)$ is the law of
  $N_j:=j+\mathrm{NB}(j+1,q)$, with
  \[
    E[N_j]=\bar n_j:=\frac{j+p}{q},\qquad
    \mathrm{Var}(N_j)=\frac{(j+1)p}{q^{2}} .
  \]
  Consequently
  $W(j,t)=\tfrac1q\,E\bigl[(1+q^{N_j})^{-1}a_{N_j}(t+L)\bigr]$.
\end{lemma}

\begin{proof}
  The substitution is valid because the window lies far above $M$;
  the swap is exact since $\chi$ has compact support, so every sum is
  finite. The indices $j\le M$ contribute super-exponentially little
  ($P(\mathrm{Bin}(n,q)=j)$ is super-exponentially small for fixed $j$
  and $n$ in the window) and are absorbed silently. The mass identity
  is the negative-binomial series
  $\sum_{k\ge0}\binom{j+k}{k}p^{k}=(1-p)^{-(j+1)}$; hence
  $q^{j+1}\binom{j+k}{k}p^{k}$, $k\ge0$, is the law of
  $\mathrm{NB}(j+1,q)$ in $k=n-j$, and its mean and variance, shifted
  by $j$, are those stated.
\end{proof}

The exact identity is the crux. The transposed weight carries mass
\emph{exactly} $1/q$ and is a genuine probability law --- being a
shifted sum of i.i.d.\ geometrics, an exceptionally well-behaved one.
The oscillatory cancellation is delivered by the recursion's own
combinatorial structure; no local limit theorem enters.

\begin{proof}[Proof of Proposition~\ref{prop:prop}]
  The transposed picture settles $\kappa$ and $r$ together: we show
  $A(t+L)=A(t)+O(e^{-t})$, and then
  $r(t)=A(t+L)-\kappa(t)A(t)=(1-\kappa(t))A(t)+O(e^{-t})=O(e^{-t})$,
  using $|\kappa-1|=O(e^{-t})$ (Lemma~\ref{lem:kappa}) and the
  boundedness of $A$.

  For real $n>0$ the map $n\mapsto a_n(t+L)$ is smooth, supported in
  the geometric window $\log n\in t+L+L\operatorname{supp}\chi$, and
  every derivative obeys $|a_n^{(k)}(t+L)|\le C_k\,n^{-k-1}$: each
  $n$-derivative produces one factor $n^{-1}$ and bounded constants
  ($\omega$, $L^{-1}$, $\|\chi^{(k)}\|_\infty$). The factor
  $(1+q^n)^{-1}=1+O(q^n)$ perturbs this super-exponentially little.

  Fix $j$ with $\bar n_j$ within $O(1)$ of the window, so $j$ is of
  order $e^{t}$. A second-order Taylor expansion of $a_n(t+L)$ about
  $n=\bar n_j$, taken under $E[\,\cdot\,]$, loses its first-order term
  because $E[N_j-\bar n_j]=0$:
  \[
    E\bigl[a_{N_j}(t+L)\bigr]=a_{\bar n_j}(t+L)
      +\tfrac12\,E\bigl[a''_{\xi_j}(t+L)\,(N_j-\bar n_j)^2\bigr],
  \]
  with $\xi_j$ between $N_j$ and $\bar n_j$, hence $\xi_j\ge j$ (both
  $N_j\ge j$ and $\bar n_j>j$). Thus $|a''_{\xi_j}|\le C_2\,j^{-3}$
  and, since $\mathrm{Var}(N_j)=O(j)$,
  \[
    \bigl|E[a_{N_j}(t+L)]-a_{\bar n_j}(t+L)\bigr|
      \le\tfrac12 C_2\,j^{-3}\,\mathrm{Var}(N_j)=O(j^{-2}).
  \]
  With the super-exponentially small $(1+q^{N_j})^{-1}$ correction,
  $W(j,t)=\tfrac1q\,a_{\bar n_j}(t+L)+O(j^{-2})$.

  It remains to undo the period shift. From $\bar n_j=(j+p)/q$,
  \[
    \log\bar n_j=\log j+L+\log(1+p/j),\qquad e^{-i\omega L}=1 ,
  \]
  and $\log(1+p/j)=O(j^{-1})$; expanding the three factors of
  $a_{\bar n_j}(t+L)$ gives
  $\tfrac1q\,a_{\bar n_j}(t+L)=a_j(t)\,(1+O(j^{-1}))=a_j(t)+O(j^{-2})$,
  the last step because $a_j(t)=O(j^{-1})$. Hence
  $W(j,t)=a_j(t)+O(j^{-2})$ uniformly in $j$, and
  \[
    A(t+L)-A(t)=\frac1L\sum_j x_j\bigl[W(j,t)-a_j(t)\bigr],
    \qquad
    |A(t+L)-A(t)|\le\frac{\|x\|_\infty}{L}\sum_j O(j^{-2}).
  \]
  The window confines $j$ to $[e^{t-cL},e^{t+cL}]$ with $c$ fixed by
  $\operatorname{supp}\chi$; all these indices exceed $e^{t-cL}$, so
  $\sum_j O(j^{-2})=O(e^{-t})$. This proves $A(t+L)=A(t)+O(e^{-t})$,
  and with it Proposition~\ref{prop:prop}.
\end{proof}

\noindent In its earlier form the remainder bound was a separate gap
--- a local central limit theorem, with one Edgeworth correction, for
the binomial log-deviation density. The transposed picture closes it:
the cancellation is supplied by the \emph{exact} negative-binomial
identity of Lemma~\ref{lem:transposed}, and the windowed coefficient
obeys the elementary propagation law $A(t+L)=A(t)+O(e^{-t})$.

\subsection{From propagation to non-convergence}

\begin{theorem}[renormalisation]\label{thm:renorm}
  Let $(x_n)$ be a bounded solution of the strategy-\All{} recursion
  for $n>M$. Then for each $t_0$ the sequence $A(t_0+kL)$ converges,
  as $k\to\infty$, to a limit $A_\infty$. Moreover: if $x_n$
  converges then $A_\infty=0$; and if $A_\infty\ne0$ then $(x_n)$ is
  not eventually monotone.
\end{theorem}

\begin{proof}
  Write $A_k=A(t_0+kL)$, $\kappa_k=\kappa(t_0+kL)$,
  $r_k=r(t_0+kL)$. By Proposition~\ref{prop:prop},
  $\sum_k|\kappa_k-1|<\infty$ and $\sum_k|r_k|<\infty$ (geometric
  series in $e^{-L}<1$). Solving $A_{k+1}=\kappa_kA_k+r_k$,
  \[
    A_k=\Bigl(\prod_{j<k}\kappa_j\Bigr)A_0
        +\sum_{j<k}\Bigl(\prod_{j<i<k}\kappa_i\Bigr)r_j
        \;\longrightarrow\;A_\infty ,
  \]
  the products and series converging absolutely. If $x_n$ converges
  then $A(t)\to0$, so $A_\infty=0$. If $A_\infty\ne0$ then
  $A_k\not\to0$, so $x_n$ does not converge; a bounded non-convergent
  sequence is not eventually monotone.
\end{proof}

%% =====================================================================
\section{The base case and the proof of
  Theorem~\ref{thm:main}}\label{sec:base}
%% =====================================================================

Theorem~\ref{thm:renorm} reduces Step~B to a single non-degeneracy:
the limit $A_\infty$ must be non-zero.

\begin{proposition}[base case]\label{prop:base}
  For every transferred sequence $(x_n)$ of Lemma~\ref{lem:transfer},
  $A_\infty\ne0$.
\end{proposition}

By linearity $A_\infty=A_\infty(b)+A_\infty(g)$, where $g:=x-b$
satisfies $g\ge0$ (the optimum dominates strategy~\All) and solves the
strategy-\All{} recursion for $n>M$. The non-negative solutions of
that recursion form a convex cone $\mathcal S$, and --- $A_\infty$
being linear --- $\mathcal C:=A_\infty(\mathcal S)\subset\mathbb C$ is
a convex cone, with $b\in\mathcal S$.

\gap{Gap D --- the non-degeneracy}{$A_\infty\ne0$ is the genuine
arithmetic heart of the proof. It has two parts. \emph{(D1)}
$A_\infty(b)\ne0$ for pure strategy~\All{} --- whence $b_{n,p}$ does
not converge and $\Phi_p$ is non-constant (Section~\ref{sec:nonconst});
equivalently $\mathcal M[\Xi](2\pi im/L)\ne0$ for some $m\ne0$.
\emph{(D2)} the cone $\mathcal C$ is \emph{pointed}:
$\mathcal C\cap(-\mathcal C)=\{0\}$. Granted both: were
$A_\infty(x)=0$, then $A_\infty(g)=-A_\infty(b)$, so $\pm A_\infty(b)$
both lie in $\mathcal C$, forcing $A_\infty(b)=0$ --- against D1.
Hence $A_\infty(x)\ne0$, which is Proposition~\ref{prop:base}. Both
D1 and D2 are open; work continues here.}

\begin{proof}[Proof of Theorem~\ref{thm:main}]
  Fix $p\in(0,\tfrac12)$. By Lemma~\ref{lem:reduction} it suffices to
  show that $(w_{n,p})$ is neither eventually non-decreasing nor
  eventually non-increasing.

  \emph{Not eventually non-decreasing} is Proposition~\ref{prop:A}.

  \emph{Not eventually non-increasing.} Suppose, for contradiction,
  $w_n\ge w_{n+1}$ for all $n\ge M$. By Lemma~\ref{lem:transfer} the
  transferred sequence $(x_n)$ satisfies $x_n=w_n$ for $n\ge M$ and
  solves the strategy-\All{} recursion for $n>M$; being a tail of a
  monotone bounded sequence, $x_n$ converges. By
  Theorem~\ref{thm:renorm} then $A_\infty=0$, contradicting
  Proposition~\ref{prop:base}. Hence no non-increasing tail exists.

  Both one-sided statements hold, so $(w_{n,p})$ is not eventually
  monotone; Lemma~\ref{lem:reduction} gives infinitely many strict
  local maxima and minima.
\end{proof}

\begin{remark}[strictness]
  Lemma~\ref{lem:reduction} delivers \emph{strict} extrema once
  $w_{n,p}\ne w_{n+1,p}$ for all large~$n$, which holds in all
  computed data and is expected for every $p\in(0,\tfrac12)$; absent
  it, Theorem~\ref{thm:main} holds with the non-strict notion of
  local extremum.
\end{remark}

%% =====================================================================
\section{Numerical evidence}\label{sec:numerics}
%% =====================================================================

The two predictions have been checked numerically in
arbitrary-precision arithmetic (\texttt{simulation/b\_asymptotics.py},
\texttt{simulation/renorm\_validation.py}). For strategy~\All{} the
log-periodic oscillation of Theorem~\ref{thm:ALL} is confirmed across
fourteen orders of magnitude in its amplitude. For the renormalisation
argument: the windowed coefficient $A(t_0+kL)$ is observed to converge
to a non-zero limit $A_\infty$ for every tested $p\in\{0.35,0.40,
0.42,0.45\}$, and $n\,|\kappa(n)-1|$ converges to the constant
$\tfrac{p}{2q}|\omega^2+i\omega|$ of Lemma~\ref{lem:kappa} --- a
direct quantitative confirmation of Proposition~\ref{prop:prop}.

%% =====================================================================
\section{Status and open problems}\label{sec:status}
%% =====================================================================

The skeleton of the proof is complete and the rigorous backbone is now
substantial: the reduction (Lemma~\ref{lem:reduction}), the elementary
Step~A (Proposition~\ref{prop:A}), the exact transfer lemma
(Lemma~\ref{lem:transfer}), the log-periodic asymptotics of
strategy~\All{} (Section~\ref{sec:ALL}: the product
identity~\eqref{eq:Bidentity} and de-Poissonisation), the factor
estimate (Lemma~\ref{lem:kappa}), the propagation estimate
(Proposition~\ref{prop:prop}, via the transposed negative-binomial
kernel), and the renormalisation theorem (Theorem~\ref{thm:renorm}).
A single gap remains open, in two parts; both are flagged in the text
and summarised here.

\begin{center}
\renewcommand{\arraystretch}{1.3}
\begin{tabular}{@{}cll@{}}
\hline
\textbf{Gap} & \textbf{What is missing} & \textbf{Where} \\
\hline
D1 & $A_\infty(b)\ne0$ (non-constancy of $\Phi_p$)
   & \S\ref{sec:ALL},\,\S\ref{sec:base} \\
D2 & the cone $\mathcal C$ is pointed & \S\ref{sec:base} \\
\hline
\end{tabular}
\end{center}

\noindent A complete proof of Theorem~\ref{thm:main} now consists of
Proposition~\ref{prop:A} (done) together with Gap~D alone. Three gaps
of earlier drafts have closed. Gap~A (a Hayman-type coefficient
transfer): the strategy-\All{} asymptotics are established in
Section~\ref{sec:ALL} by the exact product identity and analytic
de-Poissonisation, with no admissibility theory. Gap~B (a Mellin
contour shift): the same route makes it unnecessary. Gap~C (the
propagation remainder): closed by the transposed kernel
(Lemma~\ref{lem:transposed}) via the exact identity
$\sum_{n\ge j}P(\mathrm{Bin}(n,q)=j)=1/q$. The sole remaining problem
is the non-degeneracy $A_\infty\ne0$ (Gap~D), the arithmetic heart;
its two parts D1 and D2 are the current front of work.

\begin{thebibliography}{9}
\bibitem{Paper1}
  P.~Pfaffelhuber,
  \emph{Optimal strategies in the all-heads coin game},
  companion paper.
\bibitem{vanDoorn}
  E.~A.~van~Doorn,
  \emph{Optimal play of a coin-flipping game},
  preprint, arXiv:2406.14700, 2024.
\bibitem{FS}
  P.~Flajolet, R.~Sedgewick,
  \emph{Analytic Combinatorics},
  Cambridge University Press, 2009.
\bibitem{JS}
  P.~Jacquet, W.~Szpankowski,
  \emph{Analytical de-Poissonization and its applications},
  Theoret.\ Comput.\ Sci.\ \textbf{201} (1998), 1--62.
\end{thebibliography}

\end{document}
